\subsection{Shared circuits and fixed noise records}
The reproducibility archive contains the recovered original circuit generator, detector matrices, logical-action matrices, fault priors, and stored syndrome/observable samples. Every backend receives the same syndrome rows and is scored against the same full logical-output vector. The primary four-block distance-five case at $p=.006$ requests all 256 stored shots; the other original configurations request the first 64 stored shots, fixed before execution. Sample and matrix SHA-256 values are retained per run. The new scaling grid uses the same generator with independently fixed seeds: $(d,B,T)=(3,16,3),(3,32,3),(3,4,12),(5,8,10),(7,4,14)$, all at $p=.003$, with 64 shots requested per configuration. Here $T$ counts noisy slices, each with the same original transversal matching and extraction rule. A separate fixed-depth sweep adds $d=5,7,9,11$ with four blocks, three noisy slices, $p=.003$, and 64 shots per configuration, using seeds $2026110940+d$; the original distance-three instance anchors that sweep. It extends the maximum code distance to eleven without changing extraction depth along that axis.

\subsection{Exact-oracle ablation and correction selection}
Rebuilding HUF and persistent HUF are two modes of one newly implemented Python reference. They use nonnegative log-prior incidence weights, select the smallest unsatisfied Tanner component, grow every current boundary clock by the minimum residual time, merge all saturated boundaries with deterministic ties, and activate only complete internal fault columns. Both modes perform the same deterministic terminal solve. Their complete Seamwise ownership pass, fallback, radius-one repair, and likelihood filter are identical. The resulting bitstrings and merge-size histories are compared directly. Stored histories cover successfully completed solver calls, including feasible global fallbacks, rather than only accepted seam updates. Infeasible local calls rejected during initialization do not retain their intermediate histories in this prototype; their time remains included in full-call latency. This limits the completeness of the cluster-distribution audit. A zero peak in the BP-assisted variant denotes a call solved by BP before any HUF growth; the supplementary table records this early-convergence fraction. The new implementation is not the published authors' original HUF executable and does not include BP-assisted reliability updates. Consequently its performance is not a claim about an optimized belief-HUF implementation.

An additional BP-assisted residual-HUF ablation uses 30 identical native min-sum BP iterations. A converged BP correction is returned directly. Otherwise both modes subtract the same posterior hard decision from the syndrome, use absolute posterior log-odds as nonnegative growth weights, order terminal columns by that same reliability, decode the residual, and restore the hard decision. This is an explicitly specified variant, not a claim of matching every detail of the original belief-HUF implementation. Thirty-two shared circuit syndromes additionally verify bitwise agreement and syndrome consistency before the full paired pilot.

The reference uses global Python integer bitsets, explicit member scans, and diagnostic histories. It also uses minimum-label DSU representatives and re-sorts pivot keys, rather than the analyzed union-by-rank DSU and maintained sorted pivot lists. It is an auditable prototype rather than the support-local representation analyzed in Theorem~\ref{thm:work}. Wall-clock measurements therefore test the full implementation, not merely the asymptotic number of algebra operations. One hundred independently seeded feasible small Tanner graphs additionally give exact agreement of both correction vectors and complete growth histories between modes.

\subsection{External methods and applicability}
Tesseract~\cite{tesseract} is run through its native Python/C++ package on the identical undecomposed detector model, with detector beam five and priority-queue limit 200,000. It returns predictions for all monitored logical observables. Native BP+LSD2 and BP+OSD2 use the same recorded matrices and priors, 30 min-sum BP iterations, scale .625, and combination-sweep order two. The former is a 2025 clustering baseline, whereas the latter increases the OSD search order beyond the original order-zero grid. External methods decode the full detector record; HUF modes run the complete repair policy. These distinct scopes are reported explicitly rather than relabeled as equal computational workloads.

Relay-BP~\cite{relay} was also selected, but its native package could not be built in this environment because a compatible wheel and Rust toolchain were unavailable. It is recorded as unavailable, not as a failed decoder. The 2026 GARI method~\cite{gari}, ghost protocol~\cite{ghost}, and logical-product approach~\cite{fast} require verified circuit or observable adapters; they have not been benchmarked here. Merely citing these methods does not establish a performance comparison against them. The current study therefore does not complete the desired 2026 baseline coverage.

\subsection{Timing, memory, and compute-budget amendments}
Every method is run in a fresh process, sequentially. Model/solver setup is recorded separately; one warmup is excluded and every subsequent full decode duration is retained. A separate pass over the first four inputs polls whole-process RSS at 1 ms intervals, so memory polling is absent from latency measurements. RSS includes imported libraries, detector models, stored samples, and retained diagnostics; it is not a backend-only memory estimate. Peaks shorter than the polling period may be missed. Native/Python and startup-cost differences must be considered when interpreting these implementation comparisons.

The first low-noise BB144 runs had a 180-second process budget. Following slow high-noise runs, newly started children were capped at 45 seconds including startup. This resource amendment did not use their logical outcomes. Before the additional BP-assisted runs, 180 seconds was preallocated to each primary 256-shot BP-oracle variant; other new runs retain the 45-second cap. Completed fixed-sample runs are summarized with marginal 95\% Wilson intervals; incomplete runs and partial shot journals are retained separately. Because a time limit can censor difficult syndromes preferentially, completed prefixes from interrupted runs are not used as unbiased logical-error-rate evidence. The pilot uses one timing pass per backend and has neither three independent timing replicates nor controlled CPU affinity. Its p99 estimates and small-sample logical rates are descriptive measurements, not tail-latency or threshold guarantees.

The archive contains full per-shot durations, predictions, correction bitstrings, cluster histories, input hashes, source hashes, environment versions, successful summaries, and non-completion logs. The extension measures scale without assuming bounded detector degree: actual maximum row and column degrees are recorded for every completed run. Their finite values do not prove a uniform bound over a growing circuit family.


\subsection{Failure metrics across growing block counts}
The primary error metric remains failure of any monitored logical output, matching the original study. Increasing block count increases the number of monitored outputs, so those whole-circuit probabilities are not fixed-observable threshold comparisons. The supplementary CSV also records the mean logical-observable bit-error rate and monitored-output count. Correlations between observables prevent interpreting that mean as an independent-sample rate without further analysis.

The process budget includes setup, warmup, the timed sample pass, and the separate memory pass. A missing final summary can therefore result from interruption during memory instrumentation even after some or all timed inputs completed; it is not a declaration of decoding infeasibility. All per-shot journals are retained for that distinction.


\subsection{Timing follow-up and structural-work diagnostic}
After the pilot, the primary surface-code point and the high-noise BB144 point were selected for three independent-process timing repetitions, using 64 fixed inputs and randomized method order with seed 774411. The latter point illustrates the large-cluster regime observed during the pilot; this follow-up is not a preregistered survey of all configurations. The 18 repeated runs add timing replication without increasing independent logical sample counts. A separate cProfile diagnostic uses eight fixed BB144 inputs and is excluded from latency summaries. Its insertion-invocation counter quantifies repeated column processing, while its instrumented durations serve only to locate remaining implementation costs.
